\def\EditionMode{1}
\def\EditionLabel{Local math-prose proof}
\def\EditionStatus{Proposed wrappers; canonical source unchanged}
\documentclass[10pt,oneside]{book}
\usepackage{fontspec}
\setmainfont{Latin Modern Roman}
\setsansfont{Latin Modern Sans}
\usepackage{amsmath,amssymb,mathtools}
\usepackage{unicode-math}
\setmathfont{Latin Modern Math}
\usepackage{xcolor,tikz-cd,lua-ul,marginnote,pdflscape}
\usetikzlibrary{arrows.meta,decorations.pathmorphing,calc,shapes.geometric}
\tikzcdset{ampersand replacement=\&}
\tikzcdset{every label/.append style={text=black}}
\usepackage[a4paper,left=18mm,right=45mm,top=18mm,bottom=22mm,marginparwidth=28mm,marginparsep=7mm]{geometry}
\usepackage[unicode,colorlinks=true,urlcolor=black,linkcolor=black]{hyperref}
\usepackage{bookmark}
\setlength{\parindent}{0pt}
\setlength{\parskip}{4pt plus 1pt minus 1pt}
\setlength{\emergencystretch}{1.5em}
\raggedbottom
% Include final-line descenders inside the declared text-height rectangle.
\maxdepth=0pt
\definecolor{PSHand}{RGB}{25,58,132}
\definecolor{PSDelete}{RGB}{146,37,37}
\definecolor{PSTypedOld}{RGB}{95,95,95}
\ifnum\EditionMode=2\colorlet{blue}{black}\else\colorlet{blue}{PSHand}\fi
\RenewDocumentCommand{\textcircled}{m}{\tikz[baseline=(PSCircle.base)]\node[draw,circle,inner sep=.5pt,line width=.4pt](PSCircle){#1};}
\newif\ifPSInApparatus
\newcounter{psnote}
\newcounter{psvariant}
\newcounter{pslinkmark}
\newif\ifPSInLink
\newsavebox{\PSNoteBox}
\newdimen\PSNoteLimit
{\footnotesize\global\PSNoteLimit=3\baselineskip}
\newcommand{\CurrentScanURL}{}
\newcommand{\CurrentKey}{}
\newcommand{\CurrentScan}{}
\newcommand{\pssource}[3]{\psflushmargins\clearpage\gdef\CurrentKey{#1}\gdef\CurrentScan{#2}\gdef\CurrentScanURL{#3}\setcounter{psvariant}{0}\typeout{PS-SOURCE: #1}}
\newcommand{\psdivision}[2]{\pdfbookmark[0]{#1}{#2}\par\noindent{\large\bfseries #1\par}\medskip}
\newcommand{\pslinkmarker}[2]{\stepcounter{pslinkmark}\latelua{texio.write_nl('log','PS-LINK-SHIP|#1|\CurrentKey|\arabic{pslinkmark}|#2')}}
\newcommand{\pslink}[2]{\ifPSInLink\pslinkmarker{#1}{nested}#2\else\begingroup\PSInLinktrue\href{\CurrentScanURL}{\pslinkmarker{#1}{outer}#2}\endgroup\fi}
\newcommand{\scanpage}[1]{\ifnum\EditionMode=2\leavevmode\marginpar{\scriptsize\sffamily\pslink{scanpage}{scan #1}}\else\par\noindent{\small\sffamily\pslink{scanpage}{[scan #1]}}\par\medskip\fi}
\newcommand{\mspage}[1]{\ifnum\EditionMode=2\leavevmode\marginpar{\scriptsize\sffamily p.\,#1}\else{\small\sffamily[p.\,#1]}\fi}
\newcommand{\hand}[1]{{\ifnum\EditionMode=2\color{black}\else\color{PSHand}\fi #1}}
\let\over\relax
\NewDocumentCommand{\over}{O{} m m}{{\textcolor{PSTypedOld}{#2}\ensuremath{^{\!\nearrow}}\if\relax\detokenize{#1}\relax\hand{#3}\else #3\fi}}
\NewDocumentCommand{\inserted}{O{} m}{{\if\relax\detokenize{#1}\relax\hand{\ensuremath{\langle}#2\ensuremath{\rangle}}\else\ensuremath{\langle}#2\ensuremath{\rangle}\fi}}
\newcommand{\deleted}[1]{{\color{PSDelete}\ifmmode\text{[del.: }#1\text{]}\else[\textsf{del.:} #1]\fi}}
\newcommand{\xout}[1]{{\color{PSTypedOld}\ifmmode\text{[x: }#1\text{]}\else[\textsf{x:} #1]\fi}}
\let\PSMathUnderline\underline
\RenewDocumentCommand{\underline}{+m}{\ifmmode\PSMathUnderline{#1}\else\underLine{#1}\fi}
\NewDocumentCommand{\uline}{+m}{\underline{#1}}
\NewDocumentCommand{\uuline}{+m}{\ifmmode\PSMathUnderline{\PSMathUnderline{#1}}\else\underLine{\underLine{#1}}\fi}
\newunderlinetype\PSWaveLine{\leaders\hbox{\lower.7ex\hbox{\tikz[x=1pt,y=1pt,baseline=0pt]{\draw[line width=.3pt] (0,0) .. controls (1,1) and (2,1) .. (3,0) .. controls (4,-1) and (5,-1) .. (6,0);}}}}
\NewDocumentCommand{\uwave}{+m}{{\PSWaveLine #1}}
\newcommand{\unsure}[1]{\pslink{unsure}{[#1\ifmmode\,?\else\,?\fi]}}
\newcommand{\illegible}[1]{\ifmmode\text{\pslink{illegible}{[illeg.\ #1]}}\else\pslink{illegible}{[illeg.\ #1]}\fi}
\NewDocumentCommand{\diagram}{+m}{\ifmmode\vcenter{\hbox{#1}}\else\par\medskip{\centering #1\par}\medskip\psflushmargins\fi}
\newenvironment{diagramblock}{\par\begin{center}}{\end{center}\par}
\newcommand{\PSMathNotes}{}
\makeatletter
\NewDocumentCommand{\psplacednote}{m +m}{%
      \sbox{\PSNoteBox}{\parbox{\marginparwidth}{\footnotesize\raggedright\hand{\textsf{m#1.} #2}}}%
      \ifdim\dimexpr\ht\PSNoteBox+\dp\PSNoteBox\relax<\PSNoteLimit
        \leavevmode\marginpar{\usebox{\PSNoteBox}}%
      \else
        \par\nobreak\smallskip\hrule\nobreak\smallskip
        {\footnotesize\textsf{Marginal note m#1.}\quad\hand{#2}\par}
        \nobreak\smallskip\hrule\smallskip
      \fi}
\newcommand{\psflushmargins}{\ifmmode\else\ifinner\else\ifx\PSMathNotes\@empty\else\let\PSTempNotes\PSMathNotes\global\let\PSMathNotes\@empty\PSTempNotes\fi\fi\fi}
\NewDocumentCommand{\margin}{+m}{%
  \ifPSInApparatus\ifmmode\hand{\text{[marg.: }#1\text{]}}\else\hand{\textsf{[marg.: }#1\textsf{]}}\fi
  \else\stepcounter{psnote}%
    \ifmmode
      \text{\textsuperscript{m\arabic{psnote}}}%
      \begingroup\edef\PSQueue{\endgroup\noexpand\g@addto@macro\noexpand\PSMathNotes{\noexpand\psplacednote{\arabic{psnote}}{\unexpanded{$#1$}}}}\PSQueue
    \else
      \textsuperscript{m\arabic{psnote}}\psplacednote{\arabic{psnote}}{#1}%
    \fi
  \fi}
\AtBeginDocument{\everymath\expandafter{\the\everymath\aftergroup\psflushmargins}\everydisplay\expandafter{\the\everydisplay\aftergroup\psflushmargins}}
\makeatother
\newcommand{\psnumber}[2]{\expandafter\gdef\csname psnum@#1\endcsname{#2}}
\NewDocumentCommand{\var}{+m +m m m}{\ifnum\EditionMode=0 #1\else\ifnum\EditionMode=1\ifmmode\text{\textsuperscript{\tiny\csname psnum@#4\endcsname}}\else\textsuperscript{\tiny\csname psnum@#4\endcsname}\fi\fi #2\fi}
\newcommand{\psapparatusstart}{\par\bigskip\hrule\smallskip\PSInApparatustrue\begingroup\footnotesize\raggedright\textsf{Apparatus for scan \CurrentScan}\par}
\newcommand{\psapparatusend}{\par\endgroup\PSInApparatusfalse}
\NewDocumentCommand{\psentry}{m m m +m +m m}{\par\noindent\textbf{#2. #3.}\quad\textit{Typescript:} #4\quad\textit{Corrected:} \if\relax\detokenize{#5}\relax\textnormal{[omitted]}\else #5\fi\quad #6\par}
\newcommand{\psfrontmatter}{%
\hypersetup{pageanchor=false}\begin{titlepage}\raggedright
{\Large Alexander Grothendieck\par}\vspace{12mm}
{\Huge Pursuing Stacks\par}\medskip{\LARGE 1983\par}\vspace{12mm}
{\Large\EditionLabel\par}\vspace{8mm}
{\large\EditionStatus\par}\vfill
Université de Montpellier, Grothendieck archive, cotes 134-2--134-8.\par
Including the 26 source pages of R. Brown's translation at cote 134-2, sites 115--143.\par
\vspace*{2mm}\end{titlepage}\hypersetup{pageanchor=true}
\chapter*{Editorial statement}
AI-transcribed by Claude Opus 5.5 (Anthropic, effort high); checked against the scans, and this edition built, by GPT-6 Astra (OpenAI, ChatGPT Pro mode); prompts and records by Claude (Anthropic).

The witness consulted and not reproduced is the LaTeX edition of Mateo Carmona with the collaboration of Ulrik Buchholtz (arXiv 2111.01000 v2, 2021), extended from the 2015 edition by ``the scrivener'', CC0 1.0.

The translation in the appendix is R. Brown's, prepared with assistance from L. Breen and J-L. Loday, as his note says; the rights in it belong to its translator or his heirs.

The rights in the original work belong to its author or the author's representatives or heirs. Where we hold rights, in the AI-produced transcription, translation, corrections and typesetting, they are dedicated to the public domain under CC0 1.0. This edition is published for educational, non-profit purposes, and we claim no rights we do not hold. If you hold rights in this work and object to its publication here, please open a GitHub issue and we will take it down.

\section*{How to read this form}
\ifnum\EditionMode=0
The diplomatic form preserves typed and handwritten wording, cancellations, insertions, overwritten readings, underlining, doubts, and original page numbers. Blue marks handwriting; red labels handwritten cancellations; gray labels typewriter cancellations and underlying overwritten text. Angle brackets mark insertions; a superscript arrow separates an overwritten reading from its replacement. Typed insertions and typed overwriting remain plain. The disposition of long marginal notes is identified explicitly.
\else
The corrected text applies the author's insertions, replacements, cancellations, and recorded transpositions. It retains his English and notation, including consistent forms such as ``standart'', ``developped'', ``it's'' for ``its'', and ``A la''. An isolated slip is corrected only when specifically recorded and justified. Mathematical anomalies are preserved and discussed in the guide.
\ifnum\EditionMode=1
Blue continues to identify handwriting. A numbered call identifies each maximal editorial operation, with its full diplomatic and corrected readings and reason in the apparatus following the source page. Internal layers of a wholly replaced or deleted passage are preserved within that entry rather than duplicated as separate entries. These entries are generated from the same variants that govern the running text.
\else
This reading form omits the apparatus, cancelled material, superseded readings, insertion brackets, and color distinctions. It retains the resulting text, marginal notes, underlining, and unresolved doubts. Scan and source-page references appear in the margin.
\fi\fi

Every scan mark and remaining doubt links to its source page in the archive viewer. Doubts are editorial reading marks, not emendations. The source-page breaks identify the archival sequence; a source page may occupy more than one printed page. A partial working edition is explicitly labelled and contains only reviewed source pages.
\clearpage}

\begin{document}
\pssource{ps_134-6_074}{134-6/74}{https://grothendieck.umontpellier.fr/web/viewer.html?file=../134-6.pdf\#page=75}
\psnumber{ps_134-6_074-v0001}{1}
\psnumber{ps_134-6_074-v0002}{2}
\psnumber{ps_134-6_074-v0003}{3}
\psnumber{ps_134-6_074-v0004}{4}
\psnumber{ps_134-6_074-v0005}{5}
\psnumber{ps_134-6_074-v0006}{6}
\psnumber{ps_134-6_074-v0007}{7}
\psnumber{ps_134-6_074-v0008}{8}
\psnumber{ps_134-6_074-v0009}{9}
\psnumber{ps_134-6_074-v0010}{10}
\psnumber{ps_134-6_074-v0011}{11}
\psnumber{ps_134-6_074-v0012}{12}
\psnumber{ps_134-6_074-v0013}{13}
\psnumber{ps_134-6_074-v0014}{14}
\psnumber{ps_134-6_074-v0015}{15}
\psnumber{ps_134-6_074-v0016}{16}
\psnumber{ps_134-6_074-v0017}{17}
\psnumber{ps_134-6_074-v0018}{18}
% Pursuing Stacks (cotes 134-2 to 134-8): archive cote 134-6, site page 74 (PDF page 75), typescript page 392.
% typescript. Draft compared: ps5.tex:5169-5247.
% AI-transcribed by Claude Opus 5.5 (claude-opus-5-5, adaptive thinking, effort high), Message Batches API, batch ps3 (msgbatch_01XdFmuVVBAcgc7EFxdtCqsj), 2026-11-12; reviewed against both supplied scan halves; see CHECK_DECISIONS.tsv.
\scanpage{134-6/74}
\mspage{392}

(28) \[ \underline{Z}_{\hand{\mathcal{A}}} \hand{\mathbin{\tikz[baseline=-.5ex]{\node[circle,draw,inner sep=.04em] {\ensuremath{\textcolor{black}{*}}};}}} L \hand{\simeq} L \quad . \]

Using the tautological duality relation (20) between categories of
cosheaves (with values in \hand{$\mathcal{M}$}) and categories of sheaves (with values in
\hand{$\mathcal{N}=\mathcal{M}^{o}$}), we deduce accordingly an associative and unitary operation
of \hand{$\mathcal{A}$}$_{ab}$ on any category of the type \hand{$\mathcal{A}_{\mathcal{N}}$}, namely \hand{$\mathcal{N}$}-valued sheaves on \hand{$\mathcal{A}$}.
This operation is most conveniently denoted by the \uline{Hom} symbol
(29) \[ (F, K) \hand{\longmapsto} \underline{Hom}(F,K) \ : \ \hand{\mathcal{A}}_{ab}{}^{\hand{o}} \, x \, \hand{\mathcal{A}_{\mathcal{N}} \longrightarrow \mathcal{A}_{\mathcal{N}}} \quad , \]
it's explicit description in terms of K\hand{,} viewed as a functor
\[ K\colon \hand{\mathcal{A}}_{ab}{}^{o} \hand{\longrightarrow} \hand{\mathcal{N}} \]
is by
(30) \[ \underline{Hom}(F,K)\ (G) \ \hand{\cong}\var{\deleted{\xout{Hom}}}{}{deletion}{ps_134-6_074-v0001} \ K(G \otimes F\ ) \]
for K in \hand{$\mathcal{A}_{\mathcal{N}}$}, F and G in \hand{$\mathcal{A}$}$_{ab}$ . This may be written more suggestively as
(30') \[ Hom(G,\ \underline{Hom}(F,K)) \hand{\simeq} Hom(G\otimes F,\ K) \qquad , \]
similar to (22'), where we use the notation Hom (non underlined !) in
analogy to (23) for denoting K (H) , namely
(31) \[ Hom(H,K) \overset{def}{=} K(H) \qquad . \]
If confusion is feared, we may \var{\deleted{write the}}{}{deletion}{ps_134-6_074-v0002} \var{\inserted{put a}}{\hand{put a}}{insertion}{ps_134-6_074-v0003} subscript \uline{Z} in all Hom's and
\uline{Hom}'s just introduced, as well \var{\over[typed]{i}{a}}{a}{typewriter overwriting}{ps_134-6_074-v0004}s in the internal and external tensor
product operations. \hand{$\otimes$, $*$, $\circledast$}
\var{\inserted{\uline{Comments.}}}{\hand{\uline{Comments.}}}{insertion}{ps_134-6_074-v0005}
Take for instance
\[ \hand{\mathcal{N}} = (Ab) \ , \quad\text{hence}\ \hand{\mathcal{A}_{\mathcal{N}} \simeq \var{\deleted{(\unsure{A}b)}}{}{deletion}{ps_134-6_074-v0006} \mathcal{A}_{ab}} \]
by prop.1 (compare (14)). By this equivalence, the Hom in (31) is\hand{,} just
the usual Hom set corresponding to the category structure of \hand{$\mathcal{A}$}$_{ab}$ \var{\inserted{(the Hom}}{\hand{(the Hom}}{insertion}{ps_134-6_074-v0007} endowed
moreover with the structure of abelian group, coming from the fact that
\hand{$\mathcal{A}$}$_{ab}$ is \var{\xout{\illegible{1}}}{}{typewriter deletion}{ps_134-6_074-v0008} an additive category), whereas formula (30') shows that the
``external'' \uline{Hom} in this case is nothing but the \var{\xout{inter}}{}{typewriter deletion}{ps_134-6_074-v0009} usual internal \uline{Hom}
as contemplated in (24). This I hope will convince the reader of the
adequacy of the notation used\var{\over{.}{,}}{\hand{,}}{writing over}{ps_134-6_074-v0010} \hand{(}in case of exterior \var{\deleted{\xout{t\illegible{1}}}}{}{deletion}{ps_134-6_074-v0011} operation of
abelian sheaves on \hand{$\mathcal{N}$}-v\var{\over[typed]{d}{a}}{a}{typewriter overwriting}{ps_134-6_074-v0012}lued sheaves\hand{)}, and \var{\inserted{of}}{\hand{of}}{insertion}{ps_134-6_074-v0013} the convention (31). We would
like to give a similar justification for the notations (22) (23) used in
connection with operation of abelian sheaves on \hand{$\mathcal{M}$}-val\var{\over[typed]{e}{u}}{u}{typewriter overwriting}{ps_134-6_074-v0014}\var{\over[typed]{u}{e}}{e}{typewriter overwriting}{ps_134-6_074-v0015}d \hand{\uline{\textcolor{black}{co}}}sheaves,
by interpreting \hand{this} as \var{\xout{an}}{}{typewriter deletion}{ps_134-6_074-v0016} \var{\inserted[typed]{the}}{the}{typed insertion}{ps_134-6_074-v0017} internal \var{\xout{operation}}{}{typewriter deletion}{ps_134-6_074-v0018} tensor product operation in
\hand{$\mathcal{A}$}$_{ab}$ , for suitable choice of \hand{$\mathcal{N}$}. This I am affraid cannot be done for

\margin{74}

\psapparatusstart
\psentry{ps_134-6_074-v0001}{1}{deletion}{$\deleted{\xout{Hom}}$}{}{The author cancels this material by hand; it leaves the corrected running text.}
\psentry{ps_134-6_074-v0002}{2}{deletion}{\deleted{write the}}{}{The author cancels this material by hand; it leaves the corrected running text.}
\psentry{ps_134-6_074-v0003}{3}{insertion}{\inserted{put a}}{\hand{put a}}{The handwritten insertion is incorporated at its indicated position; insertion brackets leave the running text.}
\psentry{ps_134-6_074-v0004}{4}{typewriter overwriting}{\over[typed]{i}{a}}{a}{The superposed typed reading replaces the underlying reading.}
\psentry{ps_134-6_074-v0005}{5}{insertion}{\inserted{\uline{Comments.}}}{\hand{\uline{Comments.}}}{The handwritten insertion is incorporated at its indicated position; insertion brackets leave the running text.}
\psentry{ps_134-6_074-v0006}{6}{deletion}{$\deleted{(\unsure{A}b)}$}{}{The author cancels this material by hand; it leaves the corrected running text.}
\psentry{ps_134-6_074-v0007}{7}{insertion}{\inserted{(the Hom}}{\hand{(the Hom}}{The handwritten insertion is incorporated at its indicated position; insertion brackets leave the running text.}
\psentry{ps_134-6_074-v0008}{8}{typewriter deletion}{\xout{\illegible{1}}}{}{The typewriter strikes out this material; it leaves the corrected running text.}
\psentry{ps_134-6_074-v0009}{9}{typewriter deletion}{\xout{inter}}{}{The typewriter strikes out this material; it leaves the corrected running text.}
\psentry{ps_134-6_074-v0010}{10}{writing over}{\over{.}{,}}{\hand{,}}{The superposed handwritten reading replaces the underlying reading.}
\psentry{ps_134-6_074-v0011}{11}{deletion}{\deleted{\xout{t\illegible{1}}}}{}{The author cancels this material by hand; it leaves the corrected running text.}
\psentry{ps_134-6_074-v0012}{12}{typewriter overwriting}{\over[typed]{d}{a}}{a}{The superposed typed reading replaces the underlying reading.}
\psentry{ps_134-6_074-v0013}{13}{insertion}{\inserted{of}}{\hand{of}}{The handwritten insertion is incorporated at its indicated position; insertion brackets leave the running text.}
\psentry{ps_134-6_074-v0014}{14}{typewriter overwriting}{\over[typed]{e}{u}}{u}{The superposed typed reading replaces the underlying reading.}
\psentry{ps_134-6_074-v0015}{15}{typewriter overwriting}{\over[typed]{u}{e}}{e}{The superposed typed reading replaces the underlying reading.}
\psentry{ps_134-6_074-v0016}{16}{typewriter deletion}{\xout{an}}{}{The typewriter strikes out this material; it leaves the corrected running text.}
\psentry{ps_134-6_074-v0017}{17}{typed insertion}{\inserted[typed]{the}}{the}{The typed insertion is incorporated at its indicated position; insertion brackets leave the running text.}
\psentry{ps_134-6_074-v0018}{18}{typewriter deletion}{\xout{operation}}{}{The typewriter strikes out this material; it leaves the corrected running text.}
\psapparatusend

\end{document}
